\subsection{Paper 6: FCN-Based DenseNet Framework for Automated Detection and Classification of Skin Lesions in Dermoscopy Images}
\label{sec:paper_6}

\begin{tcolorbox}[colback=gray!5!white,colframe=gray!75!black,title=\textbf{Metadata \& Citation Summary}]
\textbf{Authors:} A. A. Adegun; S. Viriri \\
\textbf{Journal / Venue:} \textit{IEEE Access} (2020) \\
\textbf{DOI:} \href{https://doi.org/10.1109/ACCESS.2020.3016651}{10.1109/ACCESS.2020.3016651} \quad | \quad \textbf{IEEE Article Number:} 9167192 \\
\textbf{Thematic Category:} \textsc{Integrated Segmentation \& Dense Feature Reuse} \quad | \quad \textbf{Citation:} \cite{9167192}
\end{tcolorbox}

\subsubsection*{Research Context \& Problem Formulation}
Fuzzy lesion margins, irregular pigment borders, and occluding hair shafts degrade end-to-end classifiers that lack prior semantic lesion localization.

\paragraph{Literature Synthesis:} Automated detection of cutaneous malignancies is frequently degraded by fuzzy lesion boundaries, color variation, and artifact presence. To overcome these barriers, this research introduces an integrated two-stage deep framework combining lesion segmentation and multi-class diagnosis. The system was trained and evaluated on the benchmark ISIC 2017 (2,000 training, 150 validation, 600 test images), ISIC 2018 (2,594 training, 100 validation, 1,000 test images), and DermQuest datasets, categorizing lesions into melanoma, seborrheic keratosis, and benign nevi. Preprocessing included Dull Razor morphological artifact suppression to eliminate occluding hair shafts, followed by contrast enhancement and extensive spatial augmentations (random rotations, affine scaling). The segmentation pipeline utilizes a Fully Convolutional Network (FCN) with encoder-decoder residual skip connections, outputting binary lesion masks that crop input patches for a deep DenseNet classification model incorporating dense connectivity blocks and transition layers. The FCN segmentation achieved a 98.5\% accuracy, 92.5\% Dice coefficient, and 86.8\% Jaccard index. For classification, the framework achieved 95.0\% accuracy on ISIC 2017 and 98.5\% on ISIC 2018, with a sensitivity of 96.0\%, specificity of 97.0\%, F1-score of 95.0\%, and an AUC of 0.98. Stated limitations include high computational resource demands during dual-stage training and slight boundary over-segmentation on lesions with poorly demarcated, diffuse margins.

\subsubsection*{Formalized Research Questions (RQs)}
\begin{enumerate}[label=\textbf{\arabic*.}, leftmargin=*]
  \item \textbf{RQ1 [Integrated Two-Stage Pipeline]:} \textit{Does cascading a Fully Convolutional Network (FCN) segmentation model with a DenseNet classification network improve multi-class skin lesion recognition over single-stage models?}
  \begin{itemize}[leftmargin=1.5em]
    \item \textbf{Investigation Methodology:} Built an integrated pipeline where an FCN with encoder-decoder residual skips segments the lesion mask to crop input patches for a DenseNet classifier.
    \item \textbf{Empirical Answer \& Evidence:} Yes. The pipeline achieved state-of-the-art classification accuracy of 95.0\% on ISIC 2017 and 98.5\% on ISIC 2018, with 96.0\% sensitivity and 0.98 AUC.
  \end{itemize}
  \item \textbf{RQ2 [Dense Connectivity \& Feature Reuse]:} \textit{What is the architectural advantage of DenseNet's dense connectivity and feature reuse in capturing subtle dermoscopic biomarkers?}
  \begin{itemize}[leftmargin=1.5em]
    \item \textbf{Investigation Methodology:} Leveraged dense connectivity blocks to concatenate feature maps from all preceding layers, maximizing gradient flow and feature reuse.
    \item \textbf{Empirical Answer \& Evidence:} Dense feature reuse preserved fine textural and micro-vascular details, leading to 97.0\% specificity and 95.0\% F1-score across 3-class lesion classification.
  \end{itemize}
  \item \textbf{RQ3 [Artifact Suppression Preprocessing]:} \textit{Does DullRazor morphological hair artifact suppression improve FCN segmentation accuracy and boundary delineation?}
  \begin{itemize}[leftmargin=1.5em]
    \item \textbf{Investigation Methodology:} Applied DullRazor morphological filtering followed by contrast enhancement and affine augmentations before feeding images to the FCN.
    \item \textbf{Empirical Answer \& Evidence:} DullRazor artifact removal significantly boosted segmentation quality, yielding a 98.5\% segmentation accuracy, 92.5\% Dice coefficient, and 86.8\% Jaccard index.
  \end{itemize}
  \item \textbf{RQ4 [Computational Complexity]:} \textit{What is the computational overhead of two-stage FCN + DenseNet training, and how does it handle diffuse, poorly demarcated lesion boundaries?}
  \begin{itemize}[leftmargin=1.5em]
    \item \textbf{Investigation Methodology:} Benchmarked training memory requirements, inference latency, and segmentation error on ambiguous lesion margins.
    \item \textbf{Empirical Answer \& Evidence:} The two-stage framework demands substantial GPU memory during sequential training, and exhibits slight boundary over-segmentation on lesions with faint, diffuse margins.
  \end{itemize}
\end{enumerate}

\subsubsection*{Experimental Benchmark \& Quantitative Highlights}
\begin{itemize}[leftmargin=*]
  \item \textbf{Benchmark Repositories:} ISIC 2017 (2,750 images), ISIC 2018 (3,694 images), DermQuest
  \item \textbf{Key Empirical Results:} \textbf{Segmentation Accuracy}: 98.5\% \quad | \quad \textbf{Dice Coefficient}: 92.5\% \quad | \quad \textbf{Jaccard Index}: 86.8\% \quad | \quad \textbf{Classification Accuracy (2018)}: 98.5\% \quad | \quad \textbf{Classification Accuracy (2017)}: 95.0\% \quad | \quad \textbf{AUC}: 0.98
\end{itemize}

\subsubsection*{Open Research Questions \& Future Directions}
\begin{itemize}[leftmargin=*]
  \item \textit{Can the FCN and DenseNet stages be trained end-to-end with joint loss backpropagation to reduce training time?}
  \item \textit{How robust is the FCN segmentation when tested on low-contrast, non-dermoscopic smartphone images?}
\end{itemize}
\vspace{0.8em}\hrule\vspace{1.2em}
